\section{Proposal}

In this part I'll introduce and explain my proposal.
In the following section, I'll discuss the complete design space exploration
and why I settled on this proposal.

\subsection{Declaration and Definition of \code{operator?:}}
\begin{enumerate}
  \item \code{operator?:} must have exactly three arguments, none of them
    defaulted, and no variadic arguments.

  \item \code{operator?:} must be a non-member function (hidden friends are
    okay).

  \item At least one of the three argument types must be a user-defined type.

  \item Given a \code{operator?:(C c, T x, U y)} declaration, the expression
    \code{c ? x : y} must not match a built-in conditional operator.

  \item Given a template \code{operator?:(auto, auto, auto)} declaration, a
    specialization that matches an existing built-in conditional operator is a
    substitution failure.

  \item If the type of the first argument is \code{bool}, \code{operator?:} can
    be defaulted (\code{= default}), otherwise defaulting the operator is
    ill-formed.

  \item For a defaulted conditional operator, \code{static_cast} to the
    declared return type of both arguments (as declared) must be well-formed.
    For example, \code{int\& operator?:(bool, A\&, int\&) = default} requires that
    \code{static_cast<int\&>(declval<A\&>())} is well-formed.

  \item Explicitly deleting an \code{operator?:} overload via \code{= delete}
    is possible (as for all other function overload sets), but it is ill-formed
    to delete an existing built-in conditional operator.
\end{enumerate}

\subsection{Calling the Conditional Operator --- Overload Resolution}
\begin{enumerate}
  \item If a built-in conditional operator is valid, that operator is used and
    \emph{no name lookup} for user-defined \code{operator?:} is performed.

  \item Otherwise, perform the usual name lookup (including ADL).

  \item Add reversed candidates if the types of \code{arg2} and \code{arg3}
    differ and if applying unary \code{!} to \code{arg1} is well-formed.
    The reverse candidates use the type of \code{!arg1} rather than \code{arg1}.

  \item Perform overload resolution.
    An additional tie-break is needed if the remaining two candidates are the
    normal and reversed forms of the same \code{operator?:} and are tied under
    partial ordering.
    In that case the normal (not-reversed) form wins.

  \item If the winner is reversed the \code{c ? a : b} expression is rewritten
    to \code{!c ? b : a}.
\end{enumerate}

The first rule makes this proposal a non-breaking change.
If the \CC{}26 conditional operator applies, existing code continues to work
with the exact same semantics.
This is a reasonable approach because the existing conditional operator has
always been providing auto-generated defaulted \code{operator?:} overloads,
which we don't want to compete with:
There is a single reasonable semantic for \code{?:} and that is already
correctly implemented by the existing machinery.

If the built-in conditional operator does not apply, existing code would be
ill-formed and so there is nothing to break.
Importantly, avoiding name lookup avoids ADL, which in turn avoids template
instantiation, which can hit the corner case that could otherwise lead to
well-formed code becoming ill-formed (see \cite{P2600R0}).

\subsection{Synthesizing a Defaulted Conditional Operator}
If the expression \code{c ? <arg2-expression> : <arg3-expression>} is not a
valid built-in conditional operator expression, and lookup with subsequent
overload resolution finds a defaulted \code{operator?:} declaration with the
declared return type \type R and the declared argument types \type{bool}, \type
A, and \type B, then the expression is equivalent to:
\medskip\begin{lstlisting}
c ? [](A a) -> R { return static_cast<R>(a); }(<arg2-expression>)
  : [](B b) -> R { return static_cast<R>(b); }(<arg3-expression>)
\end{lstlisting}
which rewrites the expression to use the existing built-in conditional
operator.
Consequently, a defaulted conditional operator still implements lazy
evaluation.

\subsection{Semantic Requirement on User-Defined \code{operator?:}}
For a not-defaulted \code{operator?:(c, a, b)} function, the implementation can
assume that the result of \code{c ? a : b} and \code{!c ? b : a} are
equivalent (if the expression \code{!c} is well-formed).
A definition like
\medskip\begin{lstlisting}
R operator?:(C c, A a, B b) { return a; }
\end{lstlisting}
breaks the semantic contract of a conditional operator.
This assumption is important to enable reversed candidates.

\subsection{Calling a User-Defined Conditional Operator}
If the selected user-defined conditional operator is not defaulted, then a
normal function call is emitted.
(This implies non-lazy evaluation.)

\subsection{Taking the address of a conditional operator}
A user-defined conditional operator is still an entity like all other
user-defined operators, which can be spelled out as \code{operator?:}.
It is thus possible to take the address of a conditional operator.
Example:
\medskip\begin{lstlisting}
R operator?:(bool, T, U) = default;

void f(T t, U u) {
  R (*p1)(bool, T, U) = &operator?:; // OK
  R x = p1(true, t, u); // x = static_cast<R>(t); no lazy eval!
  R (*p2)(bool, U, T) = &operator?:; // ERROR
}
\end{lstlisting}

This is consistent with \code{operator==}:
\medskip\begin{lstlisting}
bool operator==(T, U);

void f() {
  bool (*p1)(T, U) = &operator==; // OK
  bool (*p2)(U, T) = &operator==; // ERROR
}
\end{lstlisting}

Note that it is allowed to take the address of a \emph{defaulted} \code{operator?:}.
This implies that without \code{inline}, a defaulted \code{operator?:} emits a symbol.

\subsection{Defaulted Conditional Operator Templates}
A defaulted \code{R operator?:(bool, T, U)} template always has an implicit
\medskip\begin{lstlisting}
requires is_constructible_v<R, T> && is_constructible_v<R, U>
\end{lstlisting}
constraint.
This constraint makes the template match the invariant enforced by refusing
impossible non-template defaulted condition operator declarations:
\begin{itemize}
  \item A defaulted \code{operator?:} non-template enforces the invariant ("R
    must be a genuine common type of the two branches") as a hard error at the
    declaration. Such an \code{operator?:} simply cannot exist.
  \item A defaulted \code{operator?:} template enforces the same invariant, but
    at the instantiation, as a substitution failure.
\end{itemize}

A defaulted \code{operator?:} \emph{must never exist} for impossible
conversions (template or not).
This is important for \code{std::common_type}, which is defined via \code{?:},
so it inherits whatever failure mode \code{?:} has for a non-common pair.
\begin{itemize}
  \item With the constraint: non-common pair $\Rightarrow$ SFINAE $\Rightarrow$
    \code{common_type} correctly reports "no common type."
  \item Without the constraint: non-common pair $\Rightarrow$ hard error at
    instantiation $\Rightarrow$ \code{common_type} becomes a diagnostic in the
    middle of a trait instead of the expected "not a common type."
\end{itemize}

Adding the implicit constraint is the correct choice because the operator is
defaulted and thus the body is synthesized.
The constraint matches the body, implementing a hypothetical
\code{requires(auto)}.
For not defaulted conditional operator overloads the implementation makes no
assumptions about the body of the function and therefore does not add any
implicit constraints.

% important: implicit constraint that \code{static_cast<R>(t)} and \code{static_cast<R>(u)} are valid

\subsection{Observations}
\begin{itemize}
  \item A single
    \medskip\begin{lstlisting}[style=Vc]
R operator?:(bool, T, U) = default;
    \end{lstlisting}
    declaration is sufficient to produce the same common type \code{R} for both
    argument orders \code{T, U} and \code{U, T}.

\end{itemize}
